\section{Twierdzenie Löwenheima--Skolema}\label{sec:loewenheim-skolem}
Dotychczas pytaliśmy głównie: \emph{co można udowodnić z ustalonych aksjomatów?} Teraz pytamy: \emph{jak różne światy mogą spełniać te same aksjomaty?} Twierdzenie Löwenheima--Skolema mówi, że opis pierwszego rzędu ma ograniczoną zdolność narzucania liczby elementów nieskończonego modelu. Nie jest kolejną wersją twierdzenia o niezupełności: dotyczy przede wszystkim \emph{modeli}, a nie istnienia algorytmu wypisującego dowody.

W tej i następnej sekcji dopuszczamy dowolny zbiór zdań pierwszego rzędu, niekoniecznie efektywnie aksjomatyzowalny i niekoniecznie w języku arytmetyki. Używamy zwykłej teorii mnogości z aksjomatem wyboru jako metateorii. Jest to dodatkowe tło dla wyników o mocach modeli, a nie nowe założenie wcześniejszych twierdzeń arytmetycznych.

\subsection*{Moc zbioru i podmodel elementarny: potrzebne pojęcia}\label{sec:14:1}
\begin{definicja}[Równoliczność, przeliczalność i moc]\label{res:14:1}
Dwa zbiory są \emph{równoliczne}, jeśli istnieje między nimi bijekcja: przyporządkowanie łączące każdy element jednego zbioru z dokładnie jednym elementem drugiego, bez pominięć. \emph{Moc} $|X|$ opisuje wielkość zbioru z dokładnością do takiego przyporządkowania.

Zbiór jest \emph{przeliczalnie nieskończony}, gdy jest równoliczny z $\mathbb N$; tę moc oznaczamy $\aleph_0$. Jest \emph{co najwyżej przeliczalny}, gdy jest skończony albo przeliczalnie nieskończony. Zbiór nieprzeliczalny nie daje się wypisać w jednym ciągu indeksowanym liczbami naturalnymi. Symbole $\kappa,\lambda$ będą oznaczać moce nieskończone.

Przez $|\mathcal L|$ rozumiemy liczbę symboli pozalogicznych języka. Liczba jego formuł jest nie większa niż $\max(|\mathcal L|,\aleph_0)$: każda formuła jest skończonym napisem. W szczególności język arytmetyki ma tylko przeliczalnie wiele formuł.
\end{definicja}

\begin{definicja}[Podstruktura i podmodel elementarny]\label{res:14:2}
Niech $\mathcal A,\mathcal M$ będą strukturami tego samego języka. $\mathcal A$ jest \emph{podstrukturą} $\mathcal M$, jeśli:
\begin{itemize}
\item jej niepusta dziedzina $A$ jest podzbiorem dziedziny $M$;
\item zawiera interpretacje wszystkich stałych i jest zamknięta ze względu na wszystkie funkcje struktury $\mathcal M$;
\item funkcje i relacje w $\mathcal A$ są ograniczeniami tych z $\mathcal M$.
\end{itemize}
Jest \emph{podmodelem elementarnym}, co zapisujemy $\mathcal A\preccurlyeq\mathcal M$, jeśli ponadto dla każdej formuły $\varphi(x_1,\ldots,x_n)$ i każdej krotki $\vec a$ elementów $A$ zachodzi
\[
\mathcal A\models\varphi(\vec a)\quad\Longleftrightarrow\quad\mathcal M\models\varphi(\vec a).
\]
W tym zapisie zmiennym przypisujemy elementy $\vec a$; nie zakładamy, że każdy z nich ma nazwę w wyjściowym języku. $\mathcal M$ nazywamy wtedy \emph{rozszerzeniem elementarnym} $\mathcal A$.
\end{definicja}

\begin{uwaga}[Mniej elementów, te same odpowiedzi]\label{res:14:3}
Elementarność jest znacznie silniejsza niż zachowanie samych działań. Kwantyfikatory w $\mathcal A$ przebiegają po mniejszej dziedzinie, a mimo to odpowiedzi na wszystkie pytania pierwszego rzędu z parametrami z $A$ pozostają takie same. W szczególności oba modele spełniają te same zdania bez parametrów. Sama zgodność zdań bez parametrów nie wystarcza jednak do elementarności danego włączenia podstruktur.
\end{uwaga}

\begin{lemat}[Kryterium świadków, Tarski--Vaught]\label{res:14:4}
Podstruktura $\mathcal A$ struktury $\mathcal M$ jest elementarna wtedy i tylko wtedy, gdy dla każdej formuły $\theta(y,\vec x)$ i każdej krotki $\vec a$ z $A$ zachodzi: jeśli $\mathcal M\models\exists y\,\theta(y,\vec a)$, to istnieje $b\in A$ takie, że $\mathcal M\models\theta(b,\vec a)$. Innymi słowy: każde żądanie istnienia, dotyczące parametrów już zachowanych w $A$, ma w $A$ przynajmniej jednego świadka.
\end{lemat}
\begin{proof}
Konieczność wynika z elementarności: zdanie egzystencjalne jest prawdziwe także w $\mathcal A$, więc ma tam świadka, którego własność przenosimy z powrotem do $\mathcal M$.

Dla dostateczności przeprowadzamy indukcję po budowie formuły. Najpierw indukcja po budowie termu pokazuje, że term o argumentach z $A$ ma w obu strukturach tę samą wartość: dotyczy to zmiennych i stałych, a krok funkcyjny wynika z definicji podstruktury. Formuły atomowe mają więc tę samą wartość, ponieważ relacje w $\mathcal A$ są ograniczeniami relacji z $\mathcal M$. Spójniki zachowują zgodność wartości logicznych. Jeśli $\mathcal M\models\exists y\,\theta(y,\vec a)$, kryterium daje $b\in A$ z $\mathcal M\models\theta(b,\vec a)$, a założenie indukcyjne daje $\mathcal A\models\theta(b,\vec a)$. Odwrotnie, świadka istniejącego w $\mathcal A$ przenosimy do $\mathcal M$ tym samym założeniem indukcyjnym. Kwantyfikator uniwersalny obsługuje równoważność $\forall y\,\theta\leftrightarrow\neg\exists y\,\neg\theta$.
\end{proof}

\subsection*{W dół: duży model ma mały podmodel elementarny}\label{sec:14:2}
\begin{twierdzenie}[Löwenheim--Skolem w dół]\label{res:14:5}
Niech $\mathcal M$ będzie nieskończoną strukturą języka $\mathcal L$, niech $A\subseteq M$ i niech $\kappa$ będzie mocą nieskończoną taką, że $\max(|\mathcal L|,\aleph_0,|A|)\le\kappa\le|M|$. Istnieje $\mathcal N\preccurlyeq\mathcal M$ taka, że $A\subseteq N$ oraz $|N|=\kappa$. W szczególności w co najwyżej przeliczalnym języku każdy nieskończony model ma przeliczalnie nieskończony podmodel elementarny. Można w nim dodatkowo zachować dowolny z góry wskazany co najwyżej przeliczalny zbiór elementów.
\end{twierdzenie}
\begin{proof}
Wybieramy $X_0\subseteq M$ mocy $\kappa$, zawierający $A$ i interpretacje wszystkich stałych. Następnie budujemy $X_0\subseteq X_1\subseteq X_2\subseteq\cdots$. Mając $X_n$, rozpatrujemy wszystkie formuły $\exists y\,\theta(y,\vec a)$ z parametrami z $X_n$, prawdziwe w $\mathcal M$. Dla każdej wybieramy jednego świadka w $M$ i dołączamy go do $X_{n+1}$. Nie musimy dołączać wszystkich możliwych świadków.

Przy nieskończonej mocy $\kappa$ zbiór skończonych krotek ze zbioru mocy $\kappa$ również ma moc $\kappa$. Dokładniej, $\kappa^r=\kappa$ dla każdej dodatniej skończonej arności $r$, a przeliczalna suma takich zbiorów nadal ma moc $\kappa$. Formuł jest co najwyżej $\kappa$, zatem par „formuła, krotka parametrów” jest co najwyżej $\kappa\cdot\kappa=\kappa$. Na każdym etapie dodajemy najwyżej tyle samo świadków. Indukcyjnie $|X_n|=\kappa$: nierówność w dół wynika z $X_0\subseteq X_n$, a w górę z $\kappa+\kappa=\kappa$. Dla sumy otrzymujemy $N:=\bigcup_{n\in\mathbb N}X_n$ oraz $\kappa\le|N|\le\aleph_0\cdot\kappa=\kappa$.

Używamy tu standardowej arytmetyki mocy nieskończonych z aksjomatem wyboru. Zbiór $N$ jest zamknięty ze względu na funkcje: jeśli $\vec a\in N$, to cała skończona krotka należy do pewnego $X_n$. Istotnie, dla każdego jej elementu wybieramy etap, na którym już występuje, i bierzemy największy ze skończenie wielu numerów etapów. Świadek formuły $\exists y\,(y=f(\vec a))$ należy do $X_{n+1}$. Stałe już należą do $X_0$, więc ograniczenia interpretacji dają podstrukturę $\mathcal N$. Dokładnie ten sam argument dla dowolnej $\theta$ daje kryterium z Lematu~\ref{res:14:4}. Zatem $\mathcal N\preccurlyeq\mathcal M$.
\end{proof}

\begin{uwaga}[Czy to algorytm zmniejszający model?]\label{res:14:6}
Nie. W dowodzie wybieraliśmy świadków dla formuł \emph{prawdziwych} w danym modelu. Nie zakładaliśmy procedury rozstrzygającej tę prawdziwość ani efektywnego dostępu do wszystkich jego elementów. Jest to konstrukcja teoriomnogościowa, nie obietnica programu, który z dowolnego opisu nieskończonego modelu wyprodukuje jego małą kopię. Ponadto podmodel elementarny nie musi być kopią izomorficzną całego modelu.
\end{uwaga}

\subsection*{Zwartość: od wszystkich skończonych fragmentów do modelu}\label{sec:14:3}
Do kierunku „w górę” potrzebujemy jeszcze jednego narzędzia. Podajemy jego dowód, aby nie ukrywać za słowem „zwartość” dodatkowego założenia o istnieniu modeli.

\begin{lemat}[Istnienie modelu i zwartość pierwszego rzędu]\label{res:14:7}
Każdy niesprzeczny zbiór zdań pierwszego rzędu ma model. W konsekwencji, jeżeli każdy skończony podzbiór zbioru zdań $\Gamma$ ma model, to cały $\Gamma$ ma model. Nie wymagamy rozstrzygalności ani efektywnego wyliczania $\Gamma$.
\end{lemat}
\begin{proof}
\emph{1. Dodanie nazw dla świadków.} Zaczynamy od niesprzecznego $T$ i dodajemy do języka pomocniczą stałą, aby istniał przynajmniej jeden term domknięty. Samo poszerzenie języka o nowe stałe nie tworzy sprzeczności. Gdyby pojawił się dowód sprzeczności, z niego uzyskalibyśmy dowód zdania $\exists x(x\ne x)$, w którym nie ma żadnej nowej nazwy. W skończonym dowodzie występuje tylko skończenie wiele nowych stałych. Zastępujemy je świeżymi zmiennymi nieużywanymi w tym dowodzie, z zachowaniem warunków podstawiania. Wiersze nadal są poprawne, a wniosek się nie zmienił; byłby to już dowód sprzeczności w starym języku.

Dla każdej formuły $\theta(x)$ obecnego języka, bez innych zmiennych wolnych, wprowadzamy osobną nową stałą $c_\theta$ oraz zdanie $(\exists x\,\theta(x))\to\theta(c_\theta)$. Takie zdania nazywamy \emph{aksjomatami świadków}. W kolejnym etapie robimy to samo dla formuł rozszerzonego języka, i powtarzamy procedurę dla etapów $0,1,2,\ldots$.

Wyjaśnijmy zachowanie niesprzeczności. Jeśli dodanie jednego takiego aksjomatu $\alpha=(\exists x\,\theta(x))\to\theta(c_\theta)$ do niesprzecznego $S$ prowadziłoby do sprzeczności, to Lemat~\ref{res:3:6}(c) dawałby $S\vdash\neg\alpha$. Rozwijając negację implikacji, otrzymujemy $S\vdash\neg((\exists x\,\theta(x))\to\theta(c_\theta))\leftrightarrow((\exists x\,\theta(x))\land\neg\theta(c_\theta))$. Zatem $S$ dowodzi obu członów koniunkcji. Ponieważ $c_\theta$ nie występuje w $S$ ani w $\theta(x)$, w dowodzie $\neg\theta(c_\theta)$ zastępujemy tę stałą świeżą zmienną $z$. Daje to $S\vdash\neg\theta(z)$, a generalizacja i zmiana nazwy zmiennej związanej dają $S\vdash\forall x\,\neg\theta(x)$. To przeczy już otrzymanemu $S\vdash\exists x\,\theta(x)$.

W jednym etapie stałe są różne, a formuły $\theta$ pochodzą z języka poprzedniego etapu. Dodając skończenie wiele aksjomatów po kolei, za każdym razem mamy więc stałą niewystępującą ani w poprzednich aksjomatach tego etapu, ani w dotychczasowej teorii. Skoro każde skończone rozszerzenie jest niesprzeczne, cały etap jest niesprzeczny: ewentualny dowód sprzeczności użyłby tylko skończenie wielu nowych aksjomatów. Tak samo suma wszystkich etapów $S_H$ jest niesprzeczna, bo skończony dowód mieściłby się w jednym, dostatecznie późnym etapie. Każda formuła końcowego języka pojawia się na pewnym etapie, bo jest skończonym napisem; ma więc potrzebny aksjomat świadka.

\emph{2. Wybór jednej strony każdego zdania.} Rozszerzamy $S_H$ do maksymalnego niesprzecznego zbioru zdań $\Delta$ w tym samym języku. Uzasadnia to lemat Zorna: suma łańcucha niesprzecznych rozszerzeń jest niesprzeczna, ponieważ każdy dowód używa skończenie wielu założeń, obecnych razem w jednym ogniwie łańcucha. Maksymalność oznacza, że $\Delta$ jest domknięty na konsekwencje dowodowe i dla każdego zdania zawiera dokładnie jedną z formuł $\alpha,\neg\alpha$. Rozdzielmy uzasadnienie tych dwóch własności:
\begin{itemize}
\item Jeśli $\Delta\vdash\alpha$, dołączenie $\alpha$ nie może dać sprzeczności. Każde użycie nowego aksjomatu można zastąpić jego dotychczasowym dowodem z $\Delta$. Maksymalność wymusza zatem $\alpha\in\Delta$.
\item Jeśli $\alpha\notin\Delta$, to $\Delta\cup\{\alpha\}$ jest sprzeczny, bo inaczej byłby większym niesprzecznym rozszerzeniem. Lemat~\ref{res:3:6}(c) daje $\Delta\vdash\neg\alpha$, a poprzedni punkt daje $\neg\alpha\in\Delta$. Obie strony nie mogą należeć jednocześnie, z niesprzeczności.
\end{itemize}
W szczególności $\neg\alpha\in\Delta$ wtedy i tylko wtedy, gdy $\alpha\notin\Delta$. Ponadto $\alpha\land\beta\in\Delta$ wtedy i tylko wtedy, gdy $\alpha,\beta\in\Delta$: w jedną stronę rozdzielamy koniunkcję, a w drugą ją wprowadzamy i korzystamy z domknięcia na dowody. Są to dokładnie własności potrzebne za chwilę dla spójników w lemacie o prawdzie; pozostałe spójniki wyrażamy za pomocą negacji i koniunkcji. Ten wybór nie jest na ogół efektywny.

\emph{3. Model zbudowany z nazw.} Za elementy modelu bierzemy termy domknięte, utożsamiając $t$ z $u$ dokładnie wtedy, gdy $\Delta\vdash t=u$. Klasę termu oznaczamy $[t]$. Stałą interpretujemy jako jej klasę, a pozostałe symbole przez $f([t_1],\ldots,[t_k]):=[f(t_1,\ldots,t_k)]$ oraz $R([t_1],\ldots,[t_k])\Longleftrightarrow R(t_1,\ldots,t_k)\in\Delta$.

Aksjomaty równości zapewniają niezależność od wyboru reprezentantów. Dokładniej, zwrotność, symetria i przechodniość równości czynią $t\sim u\Longleftrightarrow\Delta\vdash t=u$ relacją równoważności. Jeśli $\Delta\vdash t_i=u_i$ dla wszystkich argumentów, aksjomaty równości dają $\Delta\vdash f(\vec t)=f(\vec u)$ oraz $\Delta\vdash R(\vec t)\leftrightarrow R(\vec u)$. Pierwsza równość identyfikuje odpowiednie klasy wyników, a druga, wraz z domknięciem $\Delta$ na dowody, zachowuje przynależność formuły relacyjnej do $\Delta$. Interpretacje nie zależą więc od wybranych termów. Dziedzina jest niepusta dzięki dodanej na początku stałej.

Indukcja po budowie formuły daje \emph{lemat o prawdzie}:
\[
\mathcal M_\Delta\models\varphi([t_1],\ldots,[t_k])\quad\Longleftrightarrow\quad\varphi(t_1,\ldots,t_k)\in\Delta.
\]
Najpierw indukcja po termach pokazuje, że wartością termu $s(\vec x)$ przy argumentach $[\vec t]$ jest $[s(\vec t)]$. Dla równości atomowej korzystamy więc z definicji klas, a dla innych atomów z interpretacji relacji. Dla negacji i koniunkcji stosujemy własności przynależności do $\Delta$ udowodnione w punkcie 2. Jeżeli $\exists x\,\theta(x,\vec t)\in\Delta$, aksjomat świadka daje $\theta(c,\vec t)\in\Delta$ dla odpowiedniej stałej $c$. Ten aksjomat rzeczywiście jest w $S_H$: po podstawieniu domkniętych termów $\vec t$ formuła $\theta(x,\vec t)$ ma tylko zmienną wolną $x$ i należy do pewnego etapu języka. Jej świadek został dodany w następnym etapie. Założenie indukcyjne daje $\mathcal M_\Delta\models\theta([c],[\vec t])$, więc $[c]$ jest wymaganym elementem modelu. Odwrotnie, świadek będący klasą termu daje, przez założenie indukcyjne i wprowadzenie kwantyfikatora, zdanie egzystencjalne w $\Delta$. Kwantyfikator uniwersalny sprowadzamy do negacji i egzystencjalnego. Stąd $\mathcal M_\Delta\models T$ po pominięciu dodatkowych nazw.

\emph{4. Wniosek o zwartości.} Jeśli każdy skończony fragment $\Gamma$ ma model, $\Gamma$ jest niesprzeczny: dowód sprzeczności korzystałby ze skończonego fragmentu, a poprawność rachunku z Lematu~\ref{res:3:6}(a) wyklucza sprzeczność w jego modelu. Stosujemy teraz konstrukcję z punktów 1--3.
\end{proof}

\begin{wniosek}[Zupełność logiki nie oznacza zupełności każdej teorii]\label{res:14:8}
Dla dowolnego zbioru zdań $T$ i zdania $\varphi$ zachodzi
\[
T\vdash\varphi\quad\Longleftrightarrow\quad\text{każdy model $T$ spełnia $\varphi$}.
\]
Prawą stronę zapisujemy $T\models\varphi$.
\end{wniosek}
\begin{proof}
Kierunek w prawo jest poprawnością rachunku dowodowego. Jeśli $T\nvdash\varphi$, to $T\cup\{\neg\varphi\}$ jest niesprzeczny: inaczej twierdzenie o dedukcji i klasyczna eliminacja podwójnej negacji dałyby $T\vdash\varphi$. Lemat~\ref{res:14:7} daje więc model $T$, w którym $\varphi$ jest fałszywe.
\end{proof}

\subsection*{W górę: nieskończony model ma dowolnie duże rozszerzenia}\label{sec:14:4}
\begin{twierdzenie}[Löwenheim--Skolem w górę]\label{res:14:9}
Niech $\mathcal M$ będzie nieskończoną strukturą języka $\mathcal L$. Dla każdej mocy nieskończonej $\lambda\ge\max(|M|,|\mathcal L|,\aleph_0)$ istnieje rozszerzenie elementarne $\mathcal N$ struktury $\mathcal M$ o mocy dokładnie $\lambda$.
\end{twierdzenie}
\begin{proof}
Dodajemy stałą $c_a$ dla każdego $a\in M$. \emph{Diagram elementarny} $\mathcal M$ to zbiór wszystkich zdań $\varphi(c_{a_1},\ldots,c_{a_n})$ prawdziwych w $\mathcal M$, gdy $c_a$ nazywa $a$. Uwzględniamy wszystkie formuły, także z kwantyfikatorami i negacjami, nie tylko formuły atomowe.

Dodajemy ponadto stałe $d_i$ dla $i$ ze zbioru $I$ mocy $\lambda$ oraz zdania $d_i\ne d_j$ dla wszystkich różnych $i,j$. Każdy skończony fragment tego rozszerzonego zbioru zdań ma model: w $\mathcal M$ zachowujemy znaczenia $c_a$, a skończenie wiele występujących $d_i$ interpretujemy jako różne elementy nieskończonego $M$. Zwartość daje model $\mathcal B$ całego zbioru.

Przyporządkowanie $a\mapsto c_a^{\mathcal B}$ jest różnowartościowe, ponieważ diagram zawiera $c_a\ne c_b$ dla $a\ne b$. Oznaczmy je przez $j$. Diagram zawiera również wszystkie zdania $f(c_{a_1},\ldots,c_{a_k})=c_{f^{\mathcal M}(\vec a)}$ oraz równości z $c_{a}$ dla stałych wyjściowego języka. Zatem $j$ zachowuje funkcje i stałe, a jego obraz jest podstrukturą reduktu $\mathcal B$ do języka $\mathcal L$, czyli struktury otrzymanej przez pominięcie dodatkowych symboli, bez zmiany dziedziny. Zachowuje prawdziwość każdej formuły: jeśli formuła z parametrami jest w $\mathcal M$ fałszywa, diagram zawiera jej negację. Otrzymujemy więc elementarną kopię $\mathcal M$ wewnątrz $\mathcal B$. Elementy $d_i^{\mathcal B}$ są parami różne, zatem $|B|\ge\lambda$.

Pomijamy teraz dodatkowe symbole języka, zachowując dziedzinę $B$. Stosujemy Twierdzenie~\ref{res:14:5} do zbioru złożonego z kopii $M$ oraz wszystkich $d_i^{\mathcal B}$. Ten zbiór ma moc $\lambda$. Dostajemy podmodel $\mathcal C\preccurlyeq\mathcal B$ o mocy dokładnie $\lambda$, zawierający $j[M]$. Trzeba jeszcze sprawdzić elementarność tej kopii w $\mathcal C$, nie tylko w $\mathcal B$. Dla dowolnej formuły i parametrów z $M$ mamy
\[
\mathcal M\models\varphi(\vec a)\quad\Longleftrightarrow\quad
\mathcal B\models\varphi(j(\vec a))\quad\Longleftrightarrow\quad
\mathcal C\models\varphi(j(\vec a)).
\]
Pierwszą równoważność daje diagram, drugą elementarność $\mathcal C\preccurlyeq\mathcal B$. Zatem $j$ jest także elementarnym zanurzeniem $\mathcal M$ w $\mathcal C$.

Na końcu zamieniamy kopię na dosłowne włączenie. Wybieramy zbiór $D$ rozłączny z $M$, równoliczny z $C\setminus j[M]$, oraz bijekcję $h:M\cup D\to C$ taką, że $h(a)=j(a)$ dla $a\in M$. Przenosimy strukturę $\mathcal C$ na $M\cup D$ przez $h$:
\[
f^{\mathcal N}(\vec u)=h^{-1}(f^{\mathcal C}(h(\vec u))),\qquad
R^{\mathcal N}(\vec u)\Longleftrightarrow R^{\mathcal C}(h(\vec u)).
\]
Stałe przenosimy tak samo, za pomocą $h^{-1}$. Na elementach $M$ interpretacje zgadzają się z $\mathcal M$, bo $j$ zachowuje wszystkie symbole. Powyższa równoważność dla formuł daje $\mathcal M\preccurlyeq\mathcal N$, a bijekcja $h$ daje $|N|=\lambda$.
\end{proof}

\begin{wniosek}[Modele wszystkich mocy nieskończonych]\label{res:14:10}
Jeżeli teoria w co najwyżej przeliczalnym języku pierwszego rzędu ma choć jeden nieskończony model, to ma model każdej mocy nieskończonej.
\end{wniosek}
\begin{proof}
Najpierw wybieramy model przeliczalnie nieskończony za pomocą twierdzenia „w dół”. Następnie stosujemy twierdzenie „w górę” dla dowolnej $\lambda\ge\aleph_0$. Elementarność zachowuje wszystkie aksjomaty teorii.
\end{proof}

\begin{uwaga}[Czego wymaga ten wniosek]\label{res:14:11}
Słowo \emph{nieskończony} jest niezbędne. Istnienie modelu pięcioelementowego nie wymusza istnienia żadnego modelu nieskończonego: zdanie „istnieje dokładnie pięć elementów” można zapisać w pierwszym rzędzie. Ograniczenie na wielkość języka też ma znaczenie: nieprzeliczalnie wiele stałych, z aksjomatami stwierdzającymi, że są parami różne, wyklucza model przeliczalny.
\end{uwaga}

% Rozdziały 15--19, oryginalne strony 29--36.
